\section{Extended Discussion}
\label{app:discussion}

\subsection{Three pillars of empirical support}
\label{app:discussion:pillars}
The claims in this paper rest on three independent forms of
evidence, deliberately designed to expose disagreement.

\textbf{Theory.} Huber's 1964 asymptotic relative efficiency and the
minimax-$\delta$ relation (\eqref{eq:huber_minimax}) are theorems,
numerically verified against~\citet{huber2009robust} Table 4.1
(Appendix~\ref{app:huber_derivation}). The efficiency-robustness
tradeoff is derivable in closed form at the location-model level;
the IHDP finite-sample result amplifies but does not invalidate it.

\textbf{Synthetic stress tests.} The 17+ benchmarks referenced in
Appendix~\ref{app:experiments} probe each architectural axis
independently: contamination density, sample-size scaling, overlap,
nuisance-model choice, $\tx$ basis, prior sensitivity. They are
designed to \emph{falsify} as well as confirm; the
\texttt{normalize\_extremes} result
(\S\ref{sec:experiments:tail_signal}) is a direct falsification
outcome.

\textbf{External benchmark.} IHDP~\citep{hill2011bayesian} is the
field's standard clean-data CATE benchmark. We did not design it, we
did not tune to it, and the result reported in
\S\ref{sec:experiments:ihdp} is the first number the library
produced on the first replication without iterative adjustment.

When the three agree, the claim enters the headline tables. When
they disagree, the disagreement itself becomes a finding: the
theory-predicted $\approx 5\%$ ARE cost at $\delta = 1.345$ vs the
observed $2.2 \times$ PEHE penalty on IHDP is exactly this kind of
disagreement, and it motivated the ``structural, not tunable''
framing of \S\ref{sec:experiments:efficiency}.

\subsection{Crosswalk to repository documentation}
\label{app:discussion:crosswalk}
The library's internal \texttt{EXTENSIONS.md} predates this
paper and contains per-experiment write-ups at higher resolution
than the paper's page budget permits. For reviewer convenience:

\begin{itemize}
  \item \S\ref{sec:experiments:ihdp} $\leftrightarrow$
    \texttt{EXTENSIONS.md} \S17 (IHDP audit of Huber default)
  \item \S\ref{sec:experiments:whale} $\leftrightarrow$
    \texttt{EXTENSIONS.md} \S16 (CatBoost-Huber recovers whale)
  \item \S\ref{sec:experiments:efficiency} $\leftrightarrow$
    \texttt{EXTENSIONS.md} \S17.1--17.4 (theory and mechanism)
  \item \S\ref{sec:experiments:tail_signal} $\leftrightarrow$
    \texttt{benchmarks/results/tail\_heterogeneous\_cate.md}
    (this paper's new experiment).
\end{itemize}

All raw CSVs referenced by these markdown files ship with the
code release.

\subsection{A posteriori reflections on scope}
\label{app:discussion:scope}
Three reviewer concerns we anticipate.

\textbf{``Why not binary classification outcomes?''} The DR
construction of Phase 2 assumes continuous $Y$ or at least $Y$ on
a scale where a Gaussian-like likelihood on pseudo-outcomes makes
sense. Probit or logistic link functions can be retrofitted but
require a re-derivation of the Welsch pseudo-loss; we scoped this
out and flag it in the future-work direction on continuous
treatments (\S\ref{sec:discussion}), which shares the same
link-function structure.

\textbf{``Why not deep-learning nuisance fits?''} We ran
\texttt{MLPRegressor} and \texttt{TabTransformer}-style backbones
as part of the ablation suite (not reported in the main text for
space) and found no PEHE improvement over XGBoost or CatBoost on
tabular CATE data in the $N \le 5000$ regime tested. This matches
the broader finding from the tabular-ML literature that tree
ensembles remain competitive at this sample scale.

\textbf{``Why not a fully nonparametric $\tau(x)$?''} A GP prior
over $\tau$ is compatible with our architecture and is a natural
extension; we deliberately chose a parameterised basis to enable
interpretable subgroup contrasts and inequality probabilities.
Users who want a nonparametric $\tau$ surface can swap the basis
for a random-feature or Nystr\"om approximation; we did not
experiment with this for the current paper.
